% page 607 of the facsimile (image p-606 of the scan)
% block 162.6 x 246.3 mm ; left margin 39.4 mm
\pgc{17.780mm}{0.000mm}{
\l{\cel{50}599.}
\l{}
\l{\sou{tubero}.\cel{2}(\sou{bot.)}\cel{2}Protuberanco quan produktas, sur stipo, sur}
\l{\cel{3}trunko, la interplekturo di la fibri e la influro di la}
\l{\cel{3}tisuo celulara. - EFIS.}
\l{}
\l{\sou{tuberkloso}. (\sou{patol}.)Morbo infektiva e kontagiema, komuna a la}
\l{\cel{3}homo ed a la animali, imputata a la bacilo quan deskovris}
\l{\cel{3}Koch, ed a granulari partikulara. - DEFIRS.}
\l{}
\l{\sou{tuberkulo}.\cel{2}(\sou{bot.}) Aglomerajo pulpoza, qua konsistas (en la}
\l{\cel{3}maxim multa kazi) ye fekulo e qua kreskas, che ula planti,}
\l{\cel{3}ye la extremajo di lia radiki o rami sub-sula (kartofli+ -}
\l{\cel{3}terpomi - patati, topinamburi, e c. - II. (\sou{anat.}) Protube-}
\l{\cel{3}ranco fibroza, kartilagoze, e c., di ula parti di la}
\l{\cel{3}organismo homala. - DEFIRS.}
\l{}
\l{\sou{tuberkulino}. (\sou{medic.)}\cel{2}Extraktajo glicerinoza - facita per var-}
\l{\cel{3}migo e steriligita - de kultivaji ek tuberkloso-bacili de ra-}
\l{\cel{3}si inter-diferanta. -}
\l{}
\l{\sou{tuberoso}. (\sou{bot.}) Planto kun floro blanka, tre odoroza. L.\cel{1}\sou{poly}-}
\l{\cel{3}\sou{anthus tuberosa}. - DEFIRS.}
\l{}
\l{\sou{tubulo}. (\sou{tekn.})\cel{2}I. Tubaro qua duktasfluido a la organo qua uzos}
\l{\cel{3}olca. - II. (en\cel{1}\sou{automobilo}).\cel{1}\sou{Tubulo di alimento}\cel{1}: Tubo qua}
\l{\cel{3}duktas la mixuro exploziva, de la karburatoro tilla cilindri}
\l{\cel{3}di la motoro. - DEFS.}
\l{}
\l{\sou{tuerta}. Di qua nur un okulo esas vidiva. - S.}
\l{}
\l{\sou{tufo}. Fasko buketatra, densa, de hari, plumi, pili, flori,}
\l{\cel{3}planti, e c. - EF.}
\l{}
\l{\sou{tuko}. Peco de texuro (maxim-multa-kaze ek lino o kotono) uzata}
\l{\cel{3}heme por vishar, envelopar, kovrar, e c. - D.}
\l{}
\l{\sou{tukano}.\cel{2}(\sou{zool.})\cel{2}Grosa ucelo, en Brazilia, kun plumaro diverse-}
\l{\cel{3}kolora. - L.\cel{1}\sou{ramphastus}. - DEFIS.}
\l{}
\l{\sou{tulo}. (\sou{tekn.})\cel{2}Texuro qua konsistas ye reto tre desdensa, ek}
\l{\cel{3}fili tre tenua, de silko, lino, kotono. - DEFIRS.}
\l{}
\l{\sou{tulipo}. (\sou{bot}.)I. Planto ek la familio \textquotedbl{}liliacéi\textquotedbl{}, kun radiko}
\l{\cel{3}bulbatra, floro ovoforma, diversa-kolora. - L.\cel{1}\sou{tulipa}.-}
\l{\cel{3}II. La floro de ta planto. - DEFIRS.}
\l{}
\l{\sou{tumefaktar}. (\sou{trans.}) (\sou{patol.}) Igar organo inflesar morbale. -}
\l{\cel{3}EFIS.}
\l{}
\l{\sou{tumoro.}\cel{2}(\sou{patol.}) Influro morbala en parto di la organismo. -}
\l{\cel{3}DEFIS.}
\l{}
\l{\sou{tumulo}. I. Granda amaso artificala ek tero o petri, konoforma o}
\l{\cel{3}piramido-forma,quan onu erektis olim sur sepulteyo. - II.}
\l{\cel{3}Monteto de tero, sola, en lando plana, di qua la somito esas}
\l{\cel{3}platforma. - EFIS.}
}
